\selectlanguage{ngerman}%
\frenchspacing%
\begin{ledgroupsized}[r]{120mm}%
\footnotesize%
\pstart%
\noindent%
\textbf{Überlieferung:}%
\pend%
\end{ledgroupsized}%
\begin{ledgroupsized}[r]{114mm}%
\footnotesize%
\pstart%
\parindent -6mm%
\makebox[6mm][l]{\textit{L}}%
Aufzeichnung:
LH XXXV\;13,\;1 Bl.\;427--428.
Ein Bogen 2\textsuperscript{o};
Wasserzeichen auf Bl.\;428.
Eine Seite auf Bl.\;\href{\Scan{LH_35_13_01_427r}}{427\:r\textsuperscript{o}}.
Bl.\;\href{\Scan{LH_35_13_01_428r}}{428\:r\textsuperscript{o}} überliefert \textit{LSB} VI,\;4 N.\;357. %
Bl.\;\href{\Scan{LH_35_13_01_428r}}{428\:r\textsuperscript{o}} u.\;\href{\Scan{LH_35_13_01_428v}}{428\:v\textsuperscript{o}} eine Aufzeichnung zur Quadratur der Hyperbel, %
die voraussichtlich in einem späteren Band von \textit{LSB}\;VII veröffentlicht wird.
Bl.\;\href{\Scan{LH_35_13_01_427v}}{427\:v\textsuperscript{o}} ist unbeschrieben.
\pend%
\end{ledgroupsized}%
\vspace{4mm}%
\begin{ledgroup}%
\footnotesize%
\pstart%
\noindent%
\textbf{Datierungsgründe:}
Die Datierung der Aufzeichnung N.\;\ref{57578} beruht auf Eigenschaften des Textträgers sowie auf inhaltlichen Betrachtungen.
Siehe die ausführliche Begründung auf S.\;\pageref{57578_dat}.%
\pend%
\end{ledgroup}%
\selectlanguage{latin}%
\frenchspacing%
\vspace{6mm}%
\normalsize%
\pstart%
\noindent%
\lbrack%
\href{\Scan{LH_35_13_01_427r}}%
{427\:r\textsuperscript{o}}%
\rbrack\ %
Tensio%
\indexS{tensio repraesentans aream}
id habet notatu dignum,
quod in linea repraesentat aream,%
\indexS{area seu quadratura}
seu quadraturas,%
\indexS{quadratura}
adeo
\edtext{%
ad altiores dimensiones%
\indexS{dimensio altior}%
}{%
\lemma{ad}%
\Bfootnote{%
\textit{(1)}~altiorem dimensionem, %
\textit{(2)}~altiores dimensiones%
~\textit{L}%
}}
figuram attollit.%
\indexS{figura allata ad altiorem dimensionem}
Nam quaevis allongatio%
\indexS{allongatio particulae tensae}
(vel contractio%
\indexS{contractio particulae tensae}%
\lbrack)\rbrack\
in quavis particula%
\indexS{particulae tensa}
porrigitur in
\edlabel{LH_35_13_01_427r_jhs-1}%
sequentes.%
\edtext{}{%
\xxref{LH_35_13_01_427r_jhs-1}{LH_35_13_01_427r_jhs-2}%
\lemma{sequentes.}%
\Bfootnote{%
\lbrack/\rbrack\ %
\textit{(1)}~Ex %
\textit{(2)}~Sit filum \textit{A} %
\textit{(a)}~\textlangle ten\textrangle\ %
\textit{(b)}~\textlangle--\textrangle, sed no  %
\textit{(c)}~rectum %
\textit{(aa)}~extensum %
\textit{(bb)}~sed nondum tensum.%
~\textit{L}%
}}%
\pend%
\begin{edfigure}[2.50em][1.250em]
\hspace{-50mm}%
\includegraphics%
{gesamttex/edit_VIII-4/images/LH_35_13_01_427r_d1_427r.pdf}
\edcaption%
{\hspace{-50mm}\lbrack \textit{Fig.\:1, gestr.}\rbrack}%
\vspace{-10.375em}%
\hspace{50mm}%
\includegraphics%
{gesamttex/edit_VIII-4/images/LH_35_13_01_427r_d2_427r.pdf}
\edcaption[0.5em]{\hspace{50mm}\lbrack\textit{Fig.\:2}\rbrack}%
\end{edfigure}%
\pstart%
\noindent%
Sit filum \textit{A},
rectum
sed nondum tensum.%
\indexS{filum rectum tensum}%
\indexS{tensio fili recti}%
\edlabel{LH_35_13_01_427r_jhs-2}
Id divisum intelligatur in partes
\textit{01}, \textit{12}, \textit{23}, \textit{34}, \textit{45}.
Extendatur in \textit{B} %
ita
\edtext{%
ut
\textit{0}, \textit{1}, \textit{2} etc.
puncta extrema%
}{%
\lemma{ut}%
\Bfootnote{%
\textit{(1)}~puncta %
\bars{\textit{0}, \textit{1}, \textit{2} \textit{erg.}} extrema %
\textit{(2)}~%
\bars{\textit{0}, \textit{1}, \textit{2} etc. \textit{erg.}} puncta extrema%
~\textit{L}%
}}
partium,
loca sua servent,
atque
\edtext{%
ita
allongetur quaevis pars,%
\indexS{pars fili allongata}
ex recta facta curva;
sed nondum in rectum protendatur totum filum.%
}{%
\lemma{ita}%
\Bfootnote{%
\textit{(1)}~extendatur %
\textit{(2)}~allongetur quaevis pars, %
\textit{(a)}~sed non %
\textit{(b)}~ex recta facta curva; sed %
\textit{(aa)}~non %
\textit{(bb)}~nondum in %
\lbrack...\rbrack\ totum filum.%
~\textit{L}%
}}
Id vero jam fiat
\edtext{%
in\;\textit{C},
punctis%
\indexS{punctum transiens}
\textit{1}, \textit{2}, \textit{3}, \textit{4}, \textit{5} %
transeuntibus in
\textit{(1)}, \textit{(2)}, \textit{(3)}, \textit{(4)}, \textit{(5)}. %
Iam \textit{(0)(1)} differt ab \textit{01},
vel \textit{(1)(2)} ab \textit{12} etc.
\ledpb%
certa lege tensionis.%
\indexS{lex tensionis}%
\indexS{tensio fili recti}%
}{%
\lemma{in\;\textit{C},}%
\Bfootnote{%
\textit{(1)}~ubi %
\textit{(a)}~increm %
\textit{(b)}~\textlangle--\textrangle\ %
\textit{(2)}~punctis \textit{1}, \textit{2}, \textit{3}, \textit{4}, \textit{5} %
transeuntibus in \textit{(1)}, \textit{(2)}, \textit{(3)}, \textit{(4)}, \textit{(5)}. %
\textit{(a)}~Et \textit{01} vel \textit{02} etc. sit \textit{x}; et \textit{01} vel \textit{12} vel \textit{23} etc. erit $d\overline{x}$; \textit{(0)(1)} vel \textit{(0)(2)} vel \textit{(0)(3)} etc. sit \textit{y}; et \textit{(0)(1)} vel \textit{(1)(2)} vel \textit{(2)(3)} etc. erit $d\overline{y}$. Differentia inter $d\overline{x}$ et $d\overline{y}$ %
\textit{(aa)}~sit %
\textit{(bb)}~infinite parvas erit infi %
\textit{(b)}~Iam \textit{(0)(1)} differt ab \textit{01}, vel \textit{(1)(2)} %
\textit{(aa)}~ab %
\textit{(bb)}~a \textlangle--\textrangle\ %
\textit{(cc)}~ab \textit{12} etc. certa lege tensionis.%
~\textit{L}%
}}
At \textit{(0)(2)} ab \textit{02},
vel \textit{(0)(3)} ab \textit{03},
differt summa%
\indexS{summa discriminum}
\edtext{%
priorum discriminum.%
\indexS{discrimen longitudinum partis fili tensi}
Nam
\newline%
\hspace*{7,8mm}%
$\textit{(0)(1)}\ \text{est}\
\textit{01} + \text{discrim.}\: [\textit{01}\ \text{et}\: \textit{(0)(1)}]$;%
\edtext{}{%
\lemma{$\text{discrim.}\: [\textit{01}\; \text{et}\: \textit{(0)(1)}]$}%
\Cfootnote{%
Bei derartigen Ausdrücken stammen die eckigen Klammern von Leibniz. Wo sie versehentlich fehlen, werden sie vom Hrsg. stillschweigend ergänzt.%
}}
et
\newline%
\hspace*{7,8mm}%
$\textit{(1)(2)}\ \text{est}\
\textit{12} + \text{discrim.}\: [\textit{12}\ \text{et}\: \textit{(1)(2)}]$;%
}{%
\lemma{priorum}%
\Bfootnote{%
\textit{(1)}~differentiarum.\indexS{summa differentiarum}\indexS{differentia longitudinum partis fili tensi} %
\textit{(2)}~excessuum.\indexS{summa excessum}\indexS{excessus} %
\textit{(3)}~discriminum. Nam %
\textit{(a)}\;$\textit{(0)(1)}\ \text{est}\ \textit{01} + \text{diff.}\: \textit{(0)(1)}$ %
\textit{(b)}\;$\textit{(0)(1)}\ \text{est}\ \textit{01}\; +$ %
\textit{(aa)}~excess. %
\textit{(bb)}~$\text{discrim.}\, [\textit{01}\: \text{et} \textit{(0)(1)}]$; et %
\textit{(aaa)}\;\textit{(0)(2)}\;est %
\textit{(bbb)}\;\textit{(1)(2)}\;est %
\textit{(aaaa)}\;$\textit{02} + \text{diff.}\, [\textit{02}\: \text{et} \textit{(0)(2)}]$; %
\textit{(bbbb)}\;$\textit{12} + \text{discrim.}\, [\textit{12}\: \text{et} \textit{(1)(2)}]$;%
~\textit{L}%
}}
et
\newline%
\hspace*{7,8mm}%
$\textit{(2)(3)}\ \text{est}\
\textit{23} + \text{discrim.}\: [\textit{23}\ \text{et}\: \textit{(2)(3)}]$;%
\edtext{}{%
\lemma{$\text{est}\ \textit{23}\ +$}%
\Bfootnote{%
\textit{(1)}~diff. %
\textit{(2)}~discrim.%
~\textit{L}%
}}
et
\newline%
\hspace*{7,8mm}%
$\textit{(0)(2)}\ \text{est}\
\textit{(0)(1)} + \textit{(1)(2)}$;
et
\newline%
\hspace*{7,8mm}%
$\textit{(0)(3)}\ \text{est}\
\textit{(0)(1)} + \textit{(1)(2)} + \textit{(2)(3)}$. %
Ergo fiet
explicando%
\lbrack:\rbrack\
\newline%
\hspace*{7,8mm}%
\edtext{%
$\textit{(0)(1)} = \textit{01} + \text{discr.}\: [\textit{01}\; \text{et}\: \textit{(0)(1)}]$;%
}{%
\lemma{$\textit{(0)(1)} = \textit{01}\, +$}%
\Bfootnote{%
\textit{(1)}\;diff. %
\textit{(2)}\;$\text{discr.}\, [\textit{01}$ %
\textit{(a)}\;$+$ %
\textit{(b)}\;~$\text{et} \textit{(0)(1)}]$;%
~\textit{L}%
}}
et
\newline%
\hspace*{7,8mm}%
$\textit{(0)(2)} = \underbrace{\textit{01} + \textit{02}}_{\displaystyle\textit{02}}\: +\: \text{discr.}\: [\textit{01}\; \text{et}\: \textit{(0)(1)}] + \text{discr.}\: [\textit{12}\; \text{et}\: \textit{(1)(2)}]$;%
\edtext{}{%
\lemma{$\underbrace{\textit{01} + \textit{02}}_{\displaystyle\textit{02}}\ +$}%
\Bfootnote{%
\textit{(1)}~diff. %
\textit{(2)}~$\text{discr.}\: [\textit{01}$ %
\textit{(a)}~$+$ %
\textit{(b)}~$\text{et}\: \textit{(0)(1)}]\ +$ %
\textit{(aa)}~diff. %
\textit{(bb)}~discr.%
~\textit{L}%
}}
et
\newline%
\hspace*{7,8mm}%
$\textit{(0)(3)} = \underbrace{\textit{01} + \textit{12} + \textit{23}}_{\displaystyle\textit{03}}\: +\: \text{discr.}\: [\textit{01}\; \text{et}\: \textit{(0)(1)}]\: +\: \text{discr.}\: [\textit{12}\; \text{et}\: \textit{(1)(2)}]\: +\: \text{discr.}\: [\textit{23}\; \text{et}\: \textit{(2)(3)}]$.%
\edtext{}{%
\lemma{$\underbrace{\textit{01} + \textit{12} + \textit{23}}_{\displaystyle\textit{03}}\ +$}%
\Bfootnote{%
\textit{(1)}~diff. %
\textit{(2)}~$\text{discr.}\: [\textit{01}\; \text{et}\: \textit{(0)(1)}]\ +$ %
\textit{(a)}~diff. %
\textit{(b)}~discr.%
~\textit{L}%
}}
\newline%
Et proinde filum tensum \textit{C}
differt%
\indexS{filum rectum tensum}
\edlabel{LH_35_13_01_427r_jdf-1}%
ab eodem minus tenso \textit{A},
seu portio suae integralis \textit{(0)(3)}%
\indexS{portio integralis fili tensi}
a portione \textit{03} respondente,
\lbrack ut\rbrack\ %
summa omnium tensionum%
\indexS{summa tensionum}%
\indexS{tensio fili recti}
seu discriminum%
\indexS{summa discriminum}%
\indexS{discrimen longitudinum partis fili tensi}
partis cujusque partialis tensae%
\indexS{pars fili tensa}
a non tensa.%
\edtext{}{%
\xxref{LH_35_13_01_427r_jdf-1}{LH_35_13_01_427r_jdf-2}%
\lemma{ab}%
\Bfootnote{%
\textit{(1)}~cognomine \textit{A}\indexS{cognomen} %
\textit{(2)}~respo\textlangle ns\textrangle\ %
\textit{(3)}~\bars{%
\bars{eodem non tenso seu \textit{erg.}} %
a cognomine \textit{A} \textit{erg.}} respondente %
\textit{(a)}~\textlangle seu\textrangle\ %
\textit{(b)}~velut %
\textit{(c)}~nempe \textit{(0)(3)} ab ipso \textit{03} summa omnium discriminum %
\textit{(aa)}~int %
\textit{(bb)}~tensi\textlangle--\textrangle\ %
\textit{(cc)}~tensionum cujusque particulae %
\textit{(4)}~\bars{ab \textit{streicht Hrsg.}} eodem minus tenso \textit{A}, %
\textit{(a)}~\textlangle tota\textrangle\ portio %
\textit{(b)}~seu portio %
\textit{(aa)}~tot %
\textit{(bb)}~integral %
\textit{(cc)}~suae integralis %
\lbrack...\rbrack\ \textit{03} respondente, %
\bars{ut \textit{erg. Hrsg.}} summa omnium %
\textit{(aaa)}~discriminum seu tensionum %
\textit{(bbb)}~tensionum seu discriminum partis cujusque %
\textit{(aaaa)}~tensae %
\textit{(aaaaa)}~a non tensa %
\textit{(bbbbb)}~elementaris %
\textit{(bbbb)}~particularis %
\textit{(cccc)}~partialis tensae a non tensa. %
\textit{(aaaaa)}~\textlangle Sint jam\textrangle\ %
\textit{(bbbbb)}~Quod \textlangle si jam\textrangle\ ponamus %
\lbrack/\rbrack\ %
\textit{(ccccc)}~Sit \textit{01} et \textit{02}, \textit{x} %
\textit{(ddddd)}~\textit{01} vel \textit{02} vel \textit{03} sit \textit{x}. \textit{(0)(1)} vel \textit{(0)(2)} vel \textit{(0)(3)} sit \textit{y}. \lbrack/\rbrack\ \textit{01} vel \textit{12} vel \textit{23} erit $d\overline{x}$. \textit{(0)(1)} vel \textit{(1)(2)} vel \textit{(2)(3)} erit $d\overline{y}.$ %
\lbrack/\rbrack\ %
\textit{(eeeee)}~Quod si jam%
~\textit{L}%
}}%
\pend%
\pstart%
Quod si jam ponamus%
\edlabel{LH_35_13_01_427r_jdf-2}
particulas%
\indexS{particulae fili tensi}
seu partes partiales%
\indexS{pars partialis}%
\indexS{pars fili tensi}
esse infinitas,%
\indexS{partes infinitae comparabiles inter se}
et magnitudine%
\indexS{magnitudo partis incomparabiliter inferioris toto}
incomparabiliter inferiores toto,%
\indexS{pars incomparabiliter inferior toto}
comparabiles tamen inter se%
\indexS{partes infinitae comparabiles inter se}%
\lbrack,\rbrack\
adeoque esse elementares;%
\indexS{pars elementaris}
ponamusque praeterea
\edtext{%
tensiones%
\indexS{tensiones propinquae sibi inassignabiliter}%
\indexS{tensio fili recti}
minorem \textit{A}
et majorem \textit{C}
esse sibi inassignabiliter propinquas,%
\indexS{tensiones propinquae sibi inassignabiliter}%
}{%
\lemma{tensiones}%
\Bfootnote{%
\textit{(1)}~minorem %
\textit{(2)}~\textit{A} et \textit{C} esse sibi ina %
\textit{(3)}~minorem \textit{A} %
\lbrack...\rbrack\ inassegnabiliter propinquas,%
~\textit{L}%
}}
ita ut differentia earum%
\indexS{differentia tensionum infinitesima}
non sit
totis comparabilis,
\ledpb%
sed infinitesima;%
\indexS{differentia infinitesima}
patet
cum differentia ista%
\indexS{differentia tensionum infinitesima}
ex infinitis discriminibus%
\indexS{discrimen longitudinum partis fili tensi}
seu tensionibus ejusdem partis%
\indexS{tensio partis fili recti}
componatur,%
\indexS{differentia composita ex infinitis discriminibus}
\edtext{%
discrimina tensionum,%
\indexS{discrimen tensionum partis fili}
seu incrementum longitudinis %
particulae elementaris,%
\indexS{incrementum longitudinis particulae elementaris}%
\indexS{particulae elementares}
non esse comparabilia incremento longitudinis%
}{%
\lemma{discrimina}%
\Bfootnote{%
\textit{(1)}~esse %
\textit{(2)}~tensionum, seu %
\textit{(a)}~tensio %
\textit{(b)}~incrementum %
\textit{(aa)}~ex %
\textit{(bb)}~longitudinis %
\textit{(aaa)}~parti %
\textit{(bbb)}~elemen %
\textit{(ccc)}~particulae elementaris, %
\textit{(aaaa)}~esse incomp %
\textit{(bbbb)}~non esse comparabilia incremento longitudinis%
~\textit{L}%
}}
totius fili%
\indexS{incrementum longitudinis fili tensi}
ab initio aliquousque sumti%
\lbrack,\rbrack\
seu partis totalis%
\indexS{pars totalis}%
\lbrack;\rbrack\
vel differentiam inter \textit{23} et \textit{(2)(3)}
esse incomparabiliter minorem%
\indexS{differentia incomparabiliter minor}
differentia inter \textit{03} et \textit{(0)(3)}.%
\indexS{differentia longitudinum partis fili tensi}
\quad%
\edlabel{LH_35_13_01_427r_clcinf-1}%
Re ergo ad Calculum infinitesimalem reducta%
\indexS{calculus infinitesimalis}%
\lbrack:\rbrack%
\edlabel{LH_35_13_01_427r_clcinf-2}
\newline%
\par%
\edtabcolsep=-4pt
\begin{edtabularl}%
\hspace{1,2mm}\textit{01} &
vel %
\textit{02} \phantom{iiiiiiiiiiiiiiiiiiiiiiiiiiiiiii}&
vel %
\textit{03} &
etc. sit \textit{x}. Et
\\
\textit{(0)(1)} &
vel %
\textit{(0)(2)} &
vel %
\textit{(0)(3)} &
[etc.] sit \textit{(x)}.%
\edtext{%
}{%
\lemma{etc.}%
\Bfootnote{%
\textit{erg.\;Hrsg.}%
}}
\\
\edlabel{LH_35_13_01_427r_tab_ls-1}%
\edtext{}{%
\xxref{LH_35_13_01_427r_tab_ls-1}{LH_35_13_01_427r_tab_ls-2}%
\lemma{$\scriptstyle \textit{01} + \text{discr.}\, [\textit{01}\, \text{et}\, \textit{(0)(1)}]$ \lbrack...\rbrack\ $\scriptstyle + \text{discr.}\, [\textit{23}\, \text{et}\, \textit{(2)(3)}]$}%
\Bfootnote{%
\textit{erg.\;L}%
}}%
$\scriptstyle \textit{01} + \text{discr.}\, [\textit{01}\, \text{et}\, \textit{(0)(1)}]$ &
$\scriptstyle \textit{02} + \text{discr.}\, [\textit{01}\, \text{et}\, \textit{(0)(1)}]$ &
$\scriptstyle \text{seu}\ \textit{03} + \text{discr.}\, [\textit{01}\, \text{et}\, \textit{(0)(1)}]$ &
\\
&
$\scriptstyle \phantom{\textit{02}} + \text{discr.}\, [\textit{12}\, \text{et}\, \textit{(1)(2)}]$ &
$\scriptstyle \phantom{\text{seu}}\ \phantom{\textit{03}} + \text{discr.}\, [\textit{12}\, \text{et}\, \textit{(1)(2)}]$ &
\\
&
&
$\scriptstyle \phantom{\text{seu}}\ \phantom{\textit{03}} + \text{discr.}\, [\textit{23}\, \text{et}\, \textit{(2)(3)}]$%
\edlabel{LH_35_13_01_427r_tab_ls-2}
&
\\
\edlabel{LH_35_13_01_427r_tafel_fhlr-1}%
\edtext{}{%
\xxref{LH_35_13_01_427r_tafel_fhlr-1}{LH_35_13_01_427r_tafel_fhlr-2}%
\lemma{\textit{01} \lbrack...\rbrack\ $\scriptstyle\textit{(0)(4)} - \textit{(0)(3)}$}%
\Cfootnote{%
Vielmehr gelten folgende Äquivalenzen:
$\textit{02} - \textit{01} = \textit{12}$; %
$\textit{03} - \textit{02} = \textit{23}$; %
$\textit{04} - \textit{03} = \textit{34}$; %
$\textit{(0)(2)} - \textit{(0)(1)} = \textit{(1)(2)}$; %
$\textit{(0)(3)} - \textit{(0)(2)} = \textit{(2)(3)}$; %
$\textit{(0)(4)} - \textit{(0)(3)} = \textit{(3)(4)}$.%
}}%
\textit{01} &
vel %
\textit{12} &
vel %
\textit{23} &
etc. erit $d\overline{x}$. Et
\\
\edlabel{LH_35_13_01_427r_tab_ny-1}%
\edtext{}{%
\xxref{LH_35_13_01_427r_tab_ny-1}{LH_35_13_01_427r_tab_ny-2}%
\lemma{$\scriptstyle \text{seu}\; \textit{02} - \textit{01}$ \lbrack...\rbrack\ $\scriptstyle \text{seu}\; \textit{04} - \textit{03}$}%
\Bfootnote{%
\textit{erg.\;L}%
}}%
\!$\scriptstyle \text{seu}\; \textit{02} - \textit{01}$ &
\!$\scriptstyle \text{seu}\; \textit{03} - \textit{02}$ &
\!$\scriptstyle \text{seu}\; \textit{04} - \textit{03}$ &%
\edlabel{LH_35_13_01_427r_tab_ny-2}
\\
\textit{(0)(1)} &
vel %
\textit{(1)(2)} &
\edlabel{LH_35_13_01_427r_tab_dh-1}%
\edtext{}{%
\xxref{LH_35_13_01_427r_tab_dh-1}{LH_35_13_01_427r_tab_dh-2}%
\lemma{vel \textit{(2)(3)}}%
\Bfootnote{%
\bars{etc. \textit{erg. Hrsg.}} erit $d\overline{\textit{(x)}}$. %
\lbrack/\rbrack\ %
\bars{$\scriptstyle \text{seu}\; \textit{(0)(2)} - \textit{(0)(1)}$ %
\lbrack...\rbrack\ $\scriptstyle \text{seu}\; \textit{03} - \textit{02} + \textit{(2)(3)} - \textit{23}\ \text{seu}$ %
\bars{$\scriptstyle ...$ \textit{erg. Hrsg.}} $\scriptstyle d\,d\overline{\textit{(x)}}\: =\: d\,d\overline{x}\: +\: d\,\overline{\text{discr.\,...}}$ \textit{erg.}} %
\lbrack/\rbrack\ %
\textit{(1)}~\textit{03} vel \textit{(0)(3)} vel \textit{((0))((3))} etc. sit \textit{z}. %
\lbrack/\rbrack\ %
\textit{(2)}~\textit{x} vel \textit{(x)} vel \textit{((x))} etc. sit %
\textit{(a)}~\textit{z}
\textit{(b)}~$\overset{\scriptstyle 1}{\tilde{a}}$.%
~\textit{L}%
}}%
vel %
\textit{(2)(3)} &
[etc.] erit $d\overline{\textit{(x)}}$.%
\\
$\scriptstyle \text{seu}\; \textit{(0)(2)} - \textit{(0)(1)}$ &
$\scriptstyle \text{seu}\; \textit{(0)(3)} - \textit{(0)(2)}$ &
$\scriptstyle \text{seu}\; \textit{(0)(4)} - \textit{(0)(3)}$%
\edlabel{LH_35_13_01_427r_tafel_fhlr-2} &
$\scriptstyle \text{seu}\; d\overline{x}\: +\: \text{discr.}\: ...\ \text{Hinc}$
\\
$\scriptstyle \text{seu}\; \textit{02} - \textit{01} + \textit{(1)(2)} - \textit{12}$ %
&
$\scriptstyle \text{seu}\; \textit{03} - \textit{02} + \textit{(2)(3)} - \textit{23}$ %
&
$\scriptstyle \text{seu}\; [...]$%
&
$\scriptstyle d\,d\overline{\textit{(x)}}\: =\: d\,d\overline{x}\: +\: d\,\overline{\text{discr.\,...}}$
\\
\rule[0mm]{0mm}{5,6mm}%
\textit{x} &
vel %
\textit{(x)} &
vel %
\textit{((x))} &
etc. sit $\overset{\scriptstyle 1}{\tilde{a}}$.%
\edlabel{LH_35_13_01_427r_tab_dh-2}
\\
$\scriptstyle \text{id est}$ &
$\scriptstyle \text{id est}$ &
$\scriptstyle \text{id est}$ &
\\
$\scriptstyle +\: \textit{01}$ &
$\scriptstyle +\: \textit{(0)(1)}$ &
$\scriptstyle +\: \textit{((0))((1))}$ &
\\
$\scriptstyle +\: \textit{12}$ &
$\scriptstyle +\: \textit{(1)(2)}$ &
$\scriptstyle +\: \textit{((1))((2))}$ &
\\
$\scriptstyle +\: \textit{23}$ &
$\scriptstyle +\: \textit{(2)(3)}$ &
$\scriptstyle +\: \textit{((2))((3))}$ &
\\
$\scriptstyle \text{seu}\; \textit{03}$ &
$\scriptstyle \text{seu}\; \textit{(0)(3)}$ &
$\scriptstyle \text{seu}\; \textit{((0))((3))}$ &
\\
&
$\scriptstyle \text{seu}\ +\: \textit{03}$ &
$\scriptstyle \text{seu}\ +\: \textit{(0)(3)}$ &
\end{edtabularl}%
\noindent%
\vspace{1.5em}%
\hspace*{25mm}%
$
\scriptstyle+\;%
\left\{%
\hspace{-4mm}
\begin{array}{ll}%
&
\scriptstyle%
\text{discr.}\, [\textit{01}\: \text{et}\, \textit{(0)(1)}]
\\
&
\scriptstyle%
\text{discr.}\, [\textit{12}\: \text{et}\, \textit{(1)(2)}]
\\
&
\scriptstyle%
\text{discr.}\, [\textit{23}\: \text{et}\, \textit{(2)(3)}]
\end{array}%
\right.%
$%
\hspace*{11.5mm}%
$
\scriptstyle+\;%
\left\{%
\hspace{-4mm}
\begin{array}{ll}%
&
\scriptstyle%
\text{discr.}\, [\textit{(0)(1)}\: \text{et}\, \textit{((0))((1))}]
\\
&
\scriptstyle%
\text{discr.}\, [\textit{(1)(2)}\: \text{et}\, \textit{((1))((2))}]
\\
&
\scriptstyle%
\text{discr.}\, [\textit{(2)(3)}\: \text{et}\, \textit{((2))((3))}]
\end{array}%
\right.%
$%
\par%
\edtabcolsep=-6pt
\begin{edtabularl}%
$\textit{(x)} - x$\phantom{iiiiiiiiiiiiiiiiiiii} &
vel %
$\textit{((x))} - \textit{(x)}$\phantom{iiiiiiiiiiiiiiiiiiii} & %
\ledpb%
vel %
$\textit{(((x)))} - \textit{((x))}$ \phantom{iiiii} &
[etc.] erit $d\overline{z} = d\overset{\scriptstyle 1}{\tilde{a}}. \phantom{i} $%
\edtext{%
}{%
\lemma{etc.}%
\Bfootnote{%
\textit{erg.\;Hrsg.}%
}}
\\
$\scriptstyle \text{id est}$ %
&
$\scriptstyle \text{id est}$ %
&
&
\\
$\scriptstyle +\: \text{discr.}\, [\textit{(0)(1)}\, \text{et}\, \textit{01}]$ &
$\scriptstyle +\: \text{discr.}\, [\textit{((0))((1))}\, \text{et}\, \textit{(0)(1)}]$ &
&
\\
$\scriptstyle +\: \text{discr.}\, [\textit{(1)(2)}\, \text{et}\, \textit{12}]$ &
$\scriptstyle +\: \text{discr.}\, [\textit{((1))((2))}\, \text{et}\, \textit{(1)(2)}]$ &
&
\\
$\scriptstyle +\: \text{discr.}\, [\textit{(2)(3)}\, \text{et}\, \textit{23}]$ &
\edlabel{LH_35_13_01_427r_uxdp-1}%
$\scriptstyle +\: \text{discr.}\, [\textit{((2))((3))}\, \text{et}\, \textit{(2)(3)}]$ &
&
\end{edtabularl}%
\pend%
\vspace*{0.5em}%
\pstart%
\noindent%
\edtext{}{%
\xxref{LH_35_13_01_427r_uxdp-1}{LH_35_13_01_427r_uxdp-2}%
\lemma{$\scriptstyle +\: \text{discr.}\, [\textit{((2))((3))}\, \text{et}\, \textit{(2)(3)}]$}%
\Bfootnote{%
\lbrack/\rbrack\ %
\textit{(1)}~$\scriptstyle \textit{(0)(1)} - \textit{01}\ \text{vel}\ \textit{(0)(2)} - \textit{12}\ \text{vel}$ %
\lbrack/\rbrack\ %
\textit{(2)}~$\scriptstyle \text{discr.\ inter}$ %
\textit{(a)}~\textlangle$\scriptstyle \textit{01}$\textrangle\ %
\textit{(b)}~$\scriptstyle \textit{23}\ \text{et}$ %
\textit{(aa)}~\textlangle$\scriptstyle \textit{(0)(1)}$\textrangle\ %
\textit{(bb)}~$\scriptstyle \textit{(2)(3)}\ \text{vel\ inter}$ %
\textit{(aaa)}~\textlangle$\scriptstyle \textit{(0)(1)}$\textrangle\ %
\textit{(bbb)}~$\scriptstyle \textit{(2)(3)}\ \text{et}$ %
\textit{(aaaa)}~\textlangle$\scriptstyle \textit{((0))((1))}$\textrangle\ %
\textit{(bbbb)}~$\scriptstyle \textit{((2))((3))}$ %
\textit{(aaaaa)}~$\scriptstyle \text{etc.\ est\ discr.}$ %
\textit{(bbbbb)}~$\scriptstyle \text{vel\ inter\;\textit{((2))((3))}\ et\;\textit{(((2)))(((3)))}\ est\ discr.}$ %
\textit{(aaaaa-a)}~\textlangle$\scriptstyle \text{trans}$\textrangle\ %
\lbrack/\rbrack\ %
\textit{(bbbbb-b)}~$\scriptstyle \text{discr.\;terminorum}$%
~\textit{L}%
}}%
$%
\scriptstyle%
\text{discr.\ inter\;\textit{23}\ et\;\textit{(2)(3)}
vel\ inter\;\textit{(2)(3)}\ et\;\textit{((2))((3))}\
vel\ inter\;\textit{((2))((3))}\ et\;\textit{(((2)))(((3)))}\
est\ discr.}%
$
\pend%
\pstart%
\noindent%
$%
\scriptstyle%
\text{discr.\ terminorum
hujus seriei est discr.\;discrim.}%
$%
\edlabel{LH_35_13_01_427r_uxdp-2}
\pend%
\pstart%
\noindent%
$%
\scriptstyle%
\text{discr.\ inter\;\textit{01} et\;\textit{12}
vel inter\;\textit{12}
et\;\textit{23}
vel inter\;\textit{23} et\;\textit{34} %
est diff.\ diff.}%
$%
\pend%
\pstart%
\noindent%
$%
\scriptstyle%
\text{discr.\ terminorum hujus seriei est diff.\ diff.\ diff.}%
$%
\pend%
\vspace*{0.5em}%
\pstart%
Data lege tensionum%
\indexS{lex tensionum}
seu discriminum%
\indexS{lex discriminum}%
\lbrack,\rbrack\
summa discriminum%
\indexS{summa discriminum}
inter particulas respondentes%
\indexS{particulae fili tensi}
\edtext{%
proximorum%
\indexS{particulae respondentes filorum proximorum}
filorum%
\indexS{filum rectum tensum}
dat differentiam%
\indexS{differentia filorum}
filorum proximorum,%
\indexS{filum rectum tensum}
sed%
}{%
\lemma{proximorum filorum}%
\Bfootnote{%
\textit{(1)}~dat %
\textit{(2)}~seu %
\textit{(3)}~dat differentiam filorum
\textit{(a)}~, at %
\textit{(b)}~proximorum, sed%
~\textit{L}%
}}
si duo fila remota
ex eodem facta
sint comparanda,%
\indexS{fila comparanda}%
\indexS{comparatio filorum}
id fieri potest
summando omnia discrimina%
\indexS{summa discriminum}
proximorum filorum,
seu summam summarum%
\indexS{summa summarum discriminum}
\edtext{%
dando
discriminum%
}{%
\lemma{dando}%
\Bfootnote{%
\textit{(1)}~discrimi %
\textit{(2)}~discrimen %
\textit{(3)}~discriminum%
~\textit{L}%
}}
descendentium;%
\indexS{discrimen descendens}
\edtext{%
vel summando omnia incrementa%
\indexS{summa incrementorum}%
}{%
\lemma{vel}%
\Bfootnote{%
\textit{(1)}~summando omnia dif %
\textit{(2)}~summando omnia incrementa%
~\textit{L}%
}}
ejusdem partis,%
\indexS{incrementum longitudinis partis fili tensi}
et harum summarum omnium
rursus quaerendo summam.%
\indexS{summa summarum incrementorum}
\pend%
\pstart%
Quod si data sit linea \textit{5(5)},
vicissim hinc
\edtext{%
apparet modus%
\lbrack,\rbrack\
calculando%
\lbrack,\rbrack\
quomodo%
}{%
\lemma{apparet}%
\Bfootnote{%
\textit{(1)}~motus %
\textit{(2)}~modus %
\textit{(3)}~quomodo per calculum %
\textit{(4)}~modus %
\textit{(a)}~calculandi %
\textit{(b)}~calculando quomodo%
~\textit{L}%
}}
hac tensione%
\indexS{tensio fili recti}
nata intelligatur.
\pend%
\pstart%
Si Methodus quaeratur tangentium%
\indexS{methodus tangentium inversa}%
\indexS{tangens}
\edtext{%
inversa%
\lbrack,\rbrack\
fingend\lbrack um\rbrack\
idem filum%
\indexS{filum circumplicatum et successive tensum}%
\indexS{circumplicatio fili}%
\indexS{tensio fili circumplicati}
curvis ordinatim positione datis%
\indexS{curva data ordinatim}%
}{%
\lemma{inversa}%
\Bfootnote{%
\bars{fingendo \textit{ändert Hrsg.}} idem %
\textit{(1)}~curvae %
\textit{(2)}~filum %
\textit{(a)}~diversis %
\textit{(b)}~curvis ordinatim positione %
\textit{(aa)}~datis %
\textit{(bb)}~datis%
~\textit{L}%
}}
circumplicari et successive tendi.%
\indexS{filum circumplicatum et successive tensum}%
\indexS{circumplicatio fili}%
\indexS{tensio fili circumplicati}
Ubi investigabitur
quae ex
\edtext{%
tensioni\lbrack bu\rbrack s%
}{%
\lemma{tensionis}%
\Bfootnote{%
\textit{L\;ändert Hrsg.}%
}} %
sit necessaria%
\indexS{extensio fili circumplicati}
ut angulus directionis puncti%
\indexS{angulus directionis}
in filo tendendo
sit ad horizontem ordinatim datus
in quavis curva idem.
Hoc deinde reperto
pro curvis fili statibus%
\indexS{filum curvum tensum}
exhibebuntur recti,%
\indexS{filum rectum tensum}
ut determinatur %
quae pro
\edtext{%
quavis curva%
}{%
\lemma{quavis}%
\Bfootnote{%
\textit{(1)}~recta %
\textit{(2)}~curva%
~\textit{L}%
}}
sit longitudo fili;%
\indexS{filum tensum}%
\indexS{longitudo fili tensi}
ita demum licebit
inventam longitudinem
\edtext{%
fili,%
\indexS{filum tensum}%
\indexS{longitudo fili tensi}%
}{%
\lemma{fili}%
\Bfootnote{%
\textit{erg.\;L}%
}}
curvae respondentem%
\lbrack,\rbrack\
transferri in curvam.
\pend%
\newpage%
\vspace{3.0em}%
\pstart%
\noindent%
\centering%
*     *     *
\pend%
\vspace{4.0em}%
\selectlanguage{german}%
\frenchspacing%
\footnotesize%
\pstart%
\noindent%
\textls{Ausführliche Datierungsgründe:}%
\pagelabel{57578_dat}
Die titellose und undatierte Aufzeichnung N.\;\ref{57578} ist auf demselben Träger überliefert wie der ebenfalls undatierte Text \textit{LSB}~VI,\;4 N.\;357,\cite{01490} dessen Entstehung editorisch in die Zeitspanne \textit{1677 bis 1678} verortet wurde (mit Vorbehalt; vgl. \textit{LSB}~VI,\;4, S.\;1965.22--25).
Diese auschließlich auf inhaltliche Überlegungen gestützte Datierung ist nicht als bindend zu betrachten, und zwar aus folgenden Gründen:
\pend%
\pstart%
(1) Im Textträger der Aufzeichnung N.\;\ref{57578} ist als Wasserzeichen ein Krummstab (Bischofsstab) anzutreffen, welcher auch in dem Papier vorkommt, auf dem das eigenhändige Konzept~\textit{L} des undatierten Stücks \textit{LSB}~IV,\;7 N.\;25\cite{01491} verfasst wurde.
Dieser letztere Text ist editorisch auf die Zeit \textit{um 1697} (unter Vorbehalt) datiert worden; für diese Datierung sprechen sowohl inhaltliche Gründe wie auch das Papier, auf dem die fremdhändige Abschrift~\textit{l} überliefert ist (vgl. \textit{LSB}~IV,\;7, S.\;160.1–8).\cite{01491}
Das besagte Wasserzeichen (Krummstab) ist im Leibniz-Nachlass nach heutigem Kenntnisstand niergendwoanders nachgewiesen.
Dies lässt die Annahme zu, dass auch die Aufzeichnung N.\;\ref{57578} ungefähr im gleichen Zeitraum verfasst worden sein dürfte wie \textit{LSB}~IV,\;7 N.\;25\cite{01491}.
\pend%
\pstart%
(2) In der Aufzeichnung N.\;\ref{57578} verwendet Leibniz als technischen Ausdruck den Begriff \textit{calculus infinitesimalis} (vgl. \refpassage{LH_35_13_01_427r_clcinf-1}{LH_35_13_01_427r_clcinf-2}).
Nach heutigem Wissensstand ist diese Begrifflichkeit in Leibnizens Schrifttum nicht vor seiner Auseinandersetzung mit dem niederländischen Mathematiker Bernard Nieuwentijt\indexN{\textls{Nieuwentijt}, Bernard 1654\textendash1718} im Sommer 1695 belegt.
Nieuwentijt\indexN{\textls{Nieuwentijt}, Bernard 1654\textendash1718} hatte den Ausdruck \textit{calculus infinitesimalis} sowohl in seinen \textit{Considerationes circa analyseos ad quantitates infinite parvas applicatae principia et calculi differentialis usum in resolvendis problematibus geometricis} (Amsterdam 1964)\cite{01494} wie auch in seiner Abhandlung \textit{Analysis infinitorum seu Curvilineorum proprietates ex polygonorum natura deductae} (Amsterdam 1695)\cite{01495} explizit verwendet und in beiden Schriften kritisch auf Leibnizens Methode der Differentialrechnung Bezug genommen.
Vielleicht noch am 31.\;Mai 1695 sendete Johann Daniel Crafft\indexN{\textls{Crafft} (Craftius, Krafft, Kraftius), Johann Daniel 1624\textendash1697} aus Amsterdam\indexO{Amsterdam} Nieuwentijts\indexN{\textls{Nieuwentijt}, Bernard 1654\textendash1718} \textit{Considerationes}\cite{01494} und \textit{Analysis infinitorum}\cite{01495} an Leibniz (vgl. \textit{LSB}~III,\;6 N.\;123, S.\;383.1--2,\cite{01496} mit \textit{LSB}~III,\;6 N.\;109, S.\;339.12--15;\cite{01497} beide Abhandlungen sind als Handexemplare in der GWLB Hannover unter den Signaturen Leibn. Marg.\;149 u.\;150 aufbewahrt, wobei die \textit{Analysis infinitorum}\cite{01495} hinter dem Titelblatt Derik Makreels\indexN{\textls{Makreel}, Derik (Johannes) um 1653\textendash1717} Widmung an Leibniz im Namen des Autors mit Datum 31.\;Mai 1695 trägt).
Über Nieuwentijts\indexN{\textls{Nieuwentijt}, Bernard 1654\textendash1718} Stellungnahme berichtete Leibniz in den folgenden Wochen wiederholt:
am 11./21. Juni an Thomas Burnett of Kemney\indexN{\textls{Burnett of Kemney}, Thomas 1657\textendash1729} (\textit{LSB}I\;11 N.\;348, S.\;518.14--18);\cite{01498}
am 14./24. Juni an Antonio Magliabechi,\indexN{\textls{Magliabechi} (Magliabecchi, Magliabecki, Magliabekius, Magliabequius), Antonio 1633\textendash1714} Rudolf Christian von Bodenhausen\indexN{\textls{Bodenhausen} (Bodenus, Bodenausen), Rudolf Christian, Freiherr v. 1698} und Guillaume François de L'Hospital\indexN{\textls{L'Hospital} (L'Hôpital, Hospitalius), Guillaume François Antoine de 1661\textendash1704} (\textit{LSB}~I,\;11 N.\;350, S.\;523.5--9;\cite{01499} III,\;6 N.\;134, S.\;411.22--412.9;\cite{02500} N.\;135, S.\;415.18--416.20);\cite{02501}
am 20./30. Juni an Henri Basnage de Bauval\indexN{\textls{Basnage de Bauval} (Bellavallius, Beauval, Bauvallius), Henri 1656\textendash1710} (\textit{LSB}~II,\;3 N.\;20, S.\;47.1--5);\cite{02502}
am 21.\;Juni\,/\,1.\;Juli an Christiaan Huygens\indexN{\textls{Huygens} (Hugenius, Ugenius, Hugens, Huguens), Christiaan 1629\textendash1695} (\textit{LSB}~III,\;6 N.\;136, S.\;421.1--422.6);\cite{02503}
am 24.\;Juni\,/\,4.\;Juli an Johann Bernoulli\indexN{\textls{Bernoulli}, Johann 1667\textendash1748} (\textit{LSB}~III,\;6 N.\;137, S.\;430.14--431.7).\cite{02504}
Im Brief an Vincent Placcius\indexN{\textls{Placcius} (Plak), Vincent 1642\textendash1699} vom 25.\;Juni\,/\,5.\;Juli 1695 benannte Leibniz selbst, als er über Nieuwentijts\indexN{\textls{Nieuwentijt}, Bernard 1654\textendash1718} Abhandlungen berichtete, den \textit{calculum novum infinitesimalem vel differentialem a me inventum} (\textit{LSB}~II,\;3 N.\;21, S.\;52.11--17;\cite{01492} eine weitere explizite Verwendung des Begriffs erfolgt im Brief an Michel Angelo Fardella\indexN{\textls{Fardella}, Michel Angelo 1650\textendash1718} vom 3./13. September 1696: \textit{LSB}~II,\;3 N.\;72, S.\;193.7--10).\cite{01493}
Inzwischen war im Juniheft der \textit{Acta eruditorum} eine (von Martin Knorr\indexN{\textls{Knorr} (Knorre), Martin 1657\textendash1699} verfasste) anonyme Kurzbesprechung von Nieuwentijts\indexN{\textls{Nieuwentijt}, Bernard 1654\textendash1718} \textit{Considerationes} erschienen, in welcher explizit vom \textit{calculus infinitesimalis seu arithmetica analysisque infinitorum} die Rede war (vgl. \textit{AE} 1695, S.\;271;\cite{02505}\cite{01023} siehe hierzu Otto Menckes\indexN{\textls{Mencke} (Menken, Menkenius, Menque), Otto 1644\textendash1707} Brief an Leibniz vom 2./12. Oktober 1695, \textit{LSB}~I,\;11 N.\;498 mit Erläuterungen).\cite{02506}
Im Juliheft derselben Zeitschrift ließ dann Leibniz seine offizielle Erwiderung auf Nieuwentijts\indexN{\textls{Nieuwentijt}, Bernard 1654\textendash1718} kritische Bemerkungen veröffentlichen, welche die neue Begrifflichkeit im Titel selbst anführte: \textit{Responsio ad nonnullas difficultates a Dn. Bernardo Nieuwentijt circa methodum differentialem seu infinitesimalem motas}\ (\textit{AE} 1695, S.\;310--316, mit späteren \glqq Zusätze\grqq, \textit{AE} 1695, S\;369--372; \textit{LMG}~V, S.\;320--328).%
\pend%
\pstart%
Während Leibniz in späteren Jahren den technischen Begriff \textit{infinitesimalis} mit einer gewissen Häufigkeit benutzte, lässt sich kein früheres Vorkommnis desselben als gesichert belegen.
Dies ist etwa der Fall bei der Verwendung des Ausdrucks \textit{Analysis infinitesimalis} in einer tabellarischen Darstellung der idealen Anordnung der Bücher einer Bibliothek (\textit{LSB}~IV,\;5 N.\;81, S.\;654.6--7):\cite{02509}
Das undatierte Stück, das zu einem weiteren Komplex verwandter Aufzeichnungen und Darstellungen gehört (siehe \textit{LSB}~IV,\;5 N.\;70--81 und \textit{LSB}~IV,\;6 N.\;65--70, 72 u. 77--78), ist editorisch auf den Zeitraum \textit{um 1693} (unter Vorbehalt) datiert worden.
Diese Datierung, die auf der erschlossenen relativen Chronologie des gesamten Komplexes beruht, erweist sich aber als nicht bindend, da das Stück \textit{LSB}~IV,\;5 N.\;81 auf Papier aus der Mühle von Hehlen\indexO{Hehlen (bei Goslar)} (bei Goslar)\indexO{Goslar} überliefert ist, welches im Leibniz-Nachlass mindestens bis zum Jahr 1698 belegt ist (etwa in den drei zusammenhängenden Texten \textit{LSB}~IV,\;7 N.\;53--55\cite{02510}\cite{02511}\cite{02512} sowie beim Zeugen \textit{L\textsuperscript{1}} von \textit{LSB}~I,\;15 N.\;379\cite{02513} und beim Textträger von \textit{LSB}~I,\;15 N.\;457).\cite{02514}
Auch das Papier der Stücke \textit{LSB}~IV,\;5 N.\;76 bis N.\;79, von denen \textit{LSB}~IV,\;5 N.\;81\cite{02509} genetisch abhängt, ist bis zum Jahr 1699 belegt, wenngleich mit abnehmender Häufigkeit (siehe die Datierungsgründe von \textit{LSB}~IV,\;5 N.\;78, S.\;626.12--15).\cite{02515}
Eine Entstehung von \textit{LSB}~IV,\;5 N.\;81\cite{02509} nach dem Sommer 1695 ist demgemäß weder ausgeschlossen noch unplausibel.
\pend%
\pstart%
Aus den beiden Gründen (1) und (2) erweist sich nach jetzigem Wissensstand als unwahrscheinlich, dass die vorliegende Aufzeichnung N.\;\ref{57578} vor Juni 1695 verfasst wurde; auszuschließen ist dies dennoch nicht.
Ein genauer Terminus ante quem für die Entstehung von N.\;\ref{57578} lässt sich gegenwärtig nicht angeben.
Das oben erwähnte, seltene Wasserzeichen (Bischofsstab) im Textträger von N.\;\ref{57578} lässt jedoch die Vermutung zu, dass die Aufzeichnung vor Ende der 1690er Jahre niedergeschrieben wurde.
\pend%
\selectlanguage{latin}%
\frenchspacing%
\normalsize%
\newpage%
